% 题证摘录；来源与整理边界见相同编号分题文件。
\paragraph{Notation (Definitions 17.1, 17.4, 17.11--17.12).}
A Schwartz function on $\R$ is a complex-valued $C^\infty$ function $f$ such that $\sup_{x\in\R}|x^m f^{(n)}(x)|<\infty$ for all $m,n\in\Z_{\geq0}$. Write $\mathcal S(\R)$ for this space and
\[
\widehat f(y)=\int_\R f(x)e^{-2\pi ixy}\,dx,\qquad
\Gamma(s)=\int_0^\infty e^{-t}t^{s-1}\,dt\quad(\operatorname{Re}s>0).
\]
\begin{theorem}[17.8 (Poisson Summation Formula)]
For all $f\in\mathcal S(\R)$ we have the identity
\[
\sum_{n\in\Z}f(n)=\sum_{n\in\Z}\widehat f(n).
\]
\end{theorem}
\begin{proof}
We first note that both sums are well defined; the rapid decay property of Schwartz functions guarantees absolute convergence. Let $F(x):=\sum_{n\in\Z}f(x+n)$. Then $F$ is a periodic $C^\infty$-function, so it has a Fourier series expansion
\[
F(x)=\sum_{n\in\Z}c_ne^{2\pi inx},
\]
with Fourier coefficients
\[
c_n=\int_0^1 F(t)e^{-2\pi int}\,dt
=\int_0^1\sum_{m\in\Z}f(t+m)e^{-2\pi int}\,dt
=\int_\R f(t)e^{-2\pi int}\,dt=\widehat f(n).
\]
We then note that $\sum_{n\in\Z}f(n)=F(0)=\sum_{n\in\Z}c_n=\sum_{n\in\Z}\widehat f(n)$.
\end{proof}
\begin{lemma}[17.9]
Let $g(x):=e^{-\pi x^2}$. Then $\widehat g(y)=g(y)$.
\end{lemma}
\begin{proof}
The function $g(x)$ satisfies the first order ordinary differential equation
\[
g'+2\pi xg=0,\tag{1}
\]
with initial value $g(0)=1$. Multiplying both sides by $-i$ and taking Fourier transforms yields
\[
-i(\widehat {g'}+2\pi\widehat{xg})=-i(2\pi ix\widehat g+i\widehat g')=\widehat g'+2\pi x\widehat g=0,
\]
via Lemma 17.7. So $\widehat g$ also satisfies (1), and $\widehat g(0)=\int_\R e^{-\pi x^2}\,dx=1$, so $\widehat g=g$.
\end{proof}
\paragraph{17.1.1 Jacobi's theta function}
We now define the theta function
\[
\Theta(\tau):=\sum_{n\in\Z}e^{\pi in^2\tau}.
\]
The sum is absolutely convergent for $\operatorname{im}\tau>0$ and thus defines a holomorphic function on the upper half plane. It is easy to see that $\Theta(\tau)$ is periodic modulo 2, that is, $\Theta(\tau+2)=\Theta(\tau)$, but it it also satisfies another functional equation.
\begin{lemma}[17.10]
For all $a\in\R_{>0}$ we have $\Theta(ia)=\Theta(i/a)/\sqrt a$.
\end{lemma}
\begin{proof}
Put $g(x):=e^{-\pi x^2}$ and $h(x):=g(\sqrt a x)=e^{-\pi x^2a}$. Lemmas 17.6 and 17.9 imply
\[
\widehat h(y)=\widehat{g(\sqrt a x)}(y)=\widehat g(y/\sqrt a)/\sqrt a=g(y/\sqrt a)/\sqrt a.
\]
Plugging $\tau=ia$ into $\Theta(\tau)$ and applying Poisson summation (Theorem 17.8) yields
\[
\Theta(ia)=\sum_{n\in\Z}e^{-\pi n^2a}=\sum_{n\in\Z}h(n)
=\sum_{n\in\Z}\widehat h(n)
=\sum_{n\in\Z}g(n/\sqrt a)/\sqrt a=\Theta(i/a)/\sqrt a.
\]
\end{proof}
\paragraph{17.1.3 Completing the zeta function}
\begin{proof}[Author's derivation preceding equations (3)--(4)]
Let us now consider the function $F(s):=\pi^{-s}\Gamma(s)\zeta(2s)$, which is holomorphic on $\operatorname{Re}(s)>1/2$. In the region $\operatorname{Re}(s)>1/2$ we have an absolutely convergent sum
\begin{align*}
F(s)&=\pi^{-s}\Gamma(s)\sum_{n\geq1}n^{-2s}
=\sum_{n\geq1}(\pi n^2)^{-s}\Gamma(s)\\
&=\sum_{n\geq1}\int_0^\infty(\pi n^2)^{-s}t^{s-1}e^{-t}\,dt,
\end{align*}
and the substitution $t=\pi n^2y$ with $dt=\pi n^2dy$ yields
\[
F(s)=\sum_{n\geq1}\int_0^\infty(\pi n^2)^{-s}(\pi n^2y)^{s-1}e^{-\pi n^2y}\pi n^2\,dy
=\sum_{n\geq1}\int_0^\infty y^{s-1}e^{-\pi n^2y}\,dy.
\]
By the Fubini--Tonelli theorem, we can swap the sum and the integral to obtain
\[
F(s)=\int_0^\infty y^{s-1}\sum_{n\geq1}e^{-\pi n^2y}\,dy.
\]
We have $\Theta(iy)=\sum_{n\in\Z}e^{-\pi n^2y}=1+2\sum_{n\geq1}e^{-\pi n^2y}$, thus
\begin{align*}
F(s)&=\frac12\int_0^\infty y^{s-1}(\Theta(iy)-1)\,dy\\
&=\frac12\left(\int_0^1y^{s-1}\Theta(iy)\,dy-\frac1s+\int_1^\infty y^{s-1}(\Theta(iy)-1)\,dy\right).
\end{align*}
We now focus on the first integral on the RHS. The change of variable $t=1/y$ yields
\[
\int_0^1y^{s-1}\Theta(iy)\,dy
=\int_\infty^1 t^{1-s}\Theta(i/t)(-t^{-2})\,dt
=\int_1^\infty t^{-s-1}\Theta(i/t)\,dt.
\]
By Lemma 17.10, $\Theta(i/t)=\sqrt t\Theta(it)$, and adding $-\int_1^\infty t^{-s-1/2}\,dt+\int_1^\infty t^{-s-1/2}\,dt=0$ yields
\begin{align*}
&=\int_1^\infty t^{-s-1/2}(\Theta(it)-1)\,dt+\int_1^\infty t^{-s-1/2}\,dt\\
&=\int_1^\infty t^{-s-1/2}(\Theta(it)-1)\,dt-\frac1{1/2-s}.
\end{align*}
Plugging this back into our equation for $F(s)$ we obtain the identity
\[
F(s)=\frac12\int_1^\infty(y^{s-1}+y^{-s-1/2})(\Theta(iy)-1)\,dy-\frac1{2s}-\frac1{1-2s},
\]
valid on $\operatorname{Re}(s)>1/2$. We now observe that $F(s)=F(\frac12-s)$ for $s\ne0,\frac12$, which allows us to analytically extend $F(s)$ to a meromorphic function on $\C$ with poles only at $s=0,1/2$. Replacing $s$ with $s/2$ leads us to define the completed zeta function
\[
Z(s):=\pi^{-s/2}\Gamma(\tfrac s2)\zeta(s),\tag{3}
\]
which is meromorphic on $\C$ and satisfies the functional equation
\[
Z(s)=Z(1-s).\tag{4}
\]
It has simple poles at 0 and 1 (and no other poles).
\end{proof}